\clearpage\thispagestyle{empty}
\pdfbookmark[0]{Quellen zu Komet (November 1811)}{comet1811nov-sources}
\begin{center}\textbf{Quellen zu Komet (November 1811)}\end{center}
\noindent Diese KI-Übertragung ins heutige Deutsch folgt dem Brief von Carl Friedrich Gauß in \href{https://gdz.sub.uni-goettingen.de/id/PPN236007467}{\textit{Werke}, Band~VI} (Göttingen, 1874), S.~341--342. Die Werkausgabe gibt Göttingen, 15. November, als Briefdatum an. Der Brief steht in zwei Auszügen in Beitrag XLIX, \textit{Über den großen Cometen von 1811}, in der \textit{Monatlichen Correspondenz}, Band XXIV, November 1811, S.~\href{https://api.digitale-sammlungen.de/iiif/image/v2/bsb10538613_00521/full/1600,/0/default.jpg}{507}--\href{https://api.digitale-sammlungen.de/iiif/image/v2/bsb10538613_00531/full/1600,/0/default.jpg}{517}. Der erste Auszug beginnt auf S.~\href{https://api.digitale-sammlungen.de/iiif/image/v2/bsb10538613_00526/full/1600,/0/default.jpg}{512} und endet auf S.~\href{https://api.digitale-sammlungen.de/iiif/image/v2/bsb10538613_00527/full/1600,/0/default.jpg}{513}. Der zweite beginnt auf S.~\href{https://api.digitale-sammlungen.de/iiif/image/v2/bsb10538613_00529/full/1600,/0/default.jpg}{515} und reicht über S.~\href{https://api.digitale-sammlungen.de/iiif/image/v2/bsb10538613_00530/full/1600,/0/default.jpg}{516} bis zum ersten Absatz auf S.~517. Das digitalisierte Exemplar stammt aus der Bayerischen Staatsbibliothek.
\par\smallskip\noindent In der Zeitschrift nennt eine Fußnote zum ersten Auszug den 15. November und eine Fußnote zum zweiten den 10. und 15. November. Der Hinweis auf Burckhardt zum Kometen von 1301 ist eine redaktionelle Fußnote der Zeitschrift. Auf S.~516 steht dort der Singular „Beobachtung“, wo die Werkausgabe den Plural hat. Für die Rektaszension am 1. April gibt die Zeitschrift 119 7 32 an, die Werkausgabe 119 7 38. Die Übertragung folgt der Lesart der Werkausgabe. KI-Übertragung und KI-Satz.
\clearpage
